% TOP_PROBLEM: {"id": 30007213, "title": "Polynomial quantum synthesis of any unitary using a suitable Boolean oracle", "category_id": 16, "category": {"id": 16, "name": "physics", "display_name": "Mathematical Physics", "description": "Problems at the intersection of mathematics and physics.", "slug": "physics", "order_index": 16, "created_at": "2026-07-31T15:26:25.671Z"}, "status": "open"}
% FIRST_POSTED: 2026-10-06
% ATTEMPT_STATUS: solved
% OPENAI_CLAIM: {"schema_version":1,"source_repo":"https://github.com/openai/math","source_commit":"adc7f1241b42e322a6451854ab7e4b4c146bf78a","release_date":"2026-10-06","match":"full","resolution":"proved","independently_reviewed":false,"second_pass":"agreed","families":[{"family":"283","title":"Polynomial-time unitary synthesis from a Boolean oracle","subject":"Mathematical physics","manuscripts":[{"title":"Polynomial-Time Unitary Synthesis from a Boolean Oracle","pdf_path":"preprints/Polynomial-Time-Unitary-Synthesis-from-a-Boolean-Oracle-October-5-2026/paper.pdf","readme_path":"preprints/Polynomial-Time-Unitary-Synthesis-from-a-Boolean-Oracle-October-5-2026/README.md","date":"2026-10-05","theorem_ref":"preprints/Polynomial-Time-Unitary-Synthesis-from-a-Boolean-Oracle-October-5-2026/build/sections/01-introduction.tex:46 (thm:main)"}],"lean":null,"reasoning_trace":null}]}
\documentclass[11pt]{article}
\usepackage[margin=25mm]{geometry}
\usepackage{fontspec}
\setmainfont{Latin Modern Roman}
\usepackage{amsmath,amssymb,mathtools}
\usepackage{unicode-math}
\setmathfont{Latin Modern Math}
\usepackage{enumitem}
\usepackage{xurl,hyperref,bookmark}
\hypersetup{colorlinks=true,urlcolor=blue}
\setcounter{secnumdepth}{0}
\setlength{\parindent}{0pt}
\setlength{\parskip}{5pt}
\setlength{\emergencystretch}{3em}
\begin{document}
\section*{Polynomial quantum synthesis of any unitary using a suitable Boolean oracle}

\subsection{OpenAI's claim}
OpenAI presents a proof in the selected manuscript. A target-dependent Boolean oracle is claimed to let one uniformly generated polynomial quantum circuit approximate any $n$-qubit unitary channel to diamond error at most $1/2$.

\subsection{Outline of the argument}
Controlled tasks and programmable slots build a contracting block. Fourier fields encode a smaller-unitary representation in Boolean queries, and the reduction assembles the approximation with the fixed elementary gate set; finding the oracle efficiently is outside the theorem.

\subsection{Scope relative to this problem}
A target-dependent Boolean oracle is claimed to let one uniformly generated polynomial quantum circuit approximate any $n$-qubit unitary channel to diamond error at most $1/2$. The question retains these same parameters and conventions. Within these assumptions, this is a claimed affirmative answer to the stated question.

\subsection{Formal verification}
Family 283 has no Lean scope page in this release. Its selected result is not presented as formalized. This site has not executed the Lean code. The scope descriptions do not certify the remaining manuscript arguments.

\subsection{Status caveat}
OpenAI attributes these manuscripts to work by an internal model. Its release cautions that arguments without formalization can contain mistakes. This catalogue has not independently checked the proofs; the labels report the mathematical claims attributed to OpenAI.

\subsection{Sources}
\begin{itemize}
\item OpenAI, \emph{Polynomial-Time Unitary Synthesis from a Boolean Oracle}, 2026-10-05; citation key \texttt{OAI:Polynomial-Time-Unitary-Synthesis-from-a-Boolean-Oracle-October-5-2026}. \href{https://github.com/openai/math/blob/adc7f1241b42e322a6451854ab7e4b4c146bf78a/preprints/Polynomial-Time-Unitary-Synthesis-from-a-Boolean-Oracle-October-5-2026/paper.pdf}{PDF (GitHub)}; \href{https://github.com/openai/math/raw/adc7f1241b42e322a6451854ab7e4b4c146bf78a/preprints/Polynomial-Time-Unitary-Synthesis-from-a-Boolean-Oracle-October-5-2026/paper.pdf}{PDF (direct)}; \href{https://github.com/openai/math/blob/adc7f1241b42e322a6451854ab7e4b4c146bf78a/preprints/Polynomial-Time-Unitary-Synthesis-from-a-Boolean-Oracle-October-5-2026/README.md}{Manuscript page and citation}; \href{https://github.com/openai/math/blob/main/preprints/Polynomial-Time-Unitary-Synthesis-from-a-Boolean-Oracle-October-5-2026/paper.pdf}{Latest version}.
\item \href{https://github.com/openai/math/blob/adc7f1241b42e322a6451854ab7e4b4c146bf78a/CONTENTS.md}{OpenAI catalogue entry}.
\end{itemize}
\end{document}
